%LTeX: language=fr-FR
\documentclass[a4paper,11pt]{article}

% ---------- Préambule minimal ----------
\usepackage[french]{babel}
\usepackage{amsmath,amssymb}
\usepackage{enumitem}
\usepackage{tikz}
\usepackage{tasks}
\usetikzlibrary{arrows.meta,patterns}
\usepackage[margin=2cm]{geometry}
\usepackage{needspace}
\usepackage{currfile}
\usepackage{fancyhdr}
\usepackage{lastpage}
\usepackage{array}
\settasks{label=\alph*.}

\pagestyle{empty}
\input{\currfiledir ../def.tex}
\def \Document {Corrigé des exercices}


% Environnement "exercice" : \begin{exercice}{<numéro>} ... \end{exercice}
\newenvironment{exercice}[1]{%
  \Needspace{6\baselineskip}\par\bigskip\noindent\textbf{Exercice #1}\par\smallskip\noindent\ignorespaces
}{%
  \par
}

\newcolumntype{C}[1]{>{\centering\arraybackslash}m{#1}}

% ---------- Repères et courbes ----------
% Quadrillage + axes : \axes{xmin}{xmax}{ymin}{ymax}
\newcommand{\axes}[4]{%
  \draw[gray!35, very thin, step=1] (#1,#3) grid (#2,#4);
  \draw[-{Latex[length=2mm]}] (#1,0) -- (#2+0.4,0);
  \draw[-{Latex[length=2mm]}] (0,#3) -- (0,#4+0.4);
}

% Exercices 1 et 6 : fonction f définie sur [-7 ; 7]
\newcommand{\repereUn}[1]{%
  \begin{tikzpicture}[x=0.6cm, y=0.6cm]
    \axes{-7}{7}{-5}{5}
    \foreach \x in {-7,...,-1,1,2,...,7} \node[below] at (\x,0) {\scriptsize $\x$};
    \foreach \y in {-5,...,-1,1,2,...,5} \node[left] at (0,\y) {\scriptsize $\y$};
    \node[below left] at (0,0) {\scriptsize $0$};
    \draw[very thick] plot coordinates {(-7.00,-2.000) (-6.90,-1.809) (-6.80,-1.632) (-6.70,-1.463) (-6.60,-1.296) (-6.50,-1.125) (-6.40,-0.944) (-6.30,-0.747) (-6.20,-0.528) (-6.10,-0.281) (-6.00,0.000) (-5.90,0.307) (-5.80,0.626) (-5.70,0.950) (-5.60,1.277) (-5.50,1.600) (-5.40,1.915) (-5.30,2.218) (-5.20,2.502) (-5.10,2.765) (-5.00,3.000) (-4.90,3.213) (-4.80,3.412) (-4.70,3.598) (-4.60,3.771) (-4.50,3.931) (-4.40,4.079) (-4.30,4.214) (-4.20,4.338) (-4.10,4.449) (-4.00,4.550) (-3.90,4.640) (-3.80,4.718) (-3.70,4.787) (-3.60,4.845) (-3.50,4.894) (-3.40,4.933) (-3.30,4.963) (-3.20,4.984) (-3.10,4.996) (-3.00,5.000) (-2.90,4.957) (-2.80,4.840) (-2.70,4.662) (-2.60,4.440) (-2.50,4.188) (-2.40,3.920) (-2.30,3.653) (-2.20,3.400) (-2.10,3.178) (-2.00,3.000) (-1.90,2.850) (-1.80,2.704) (-1.70,2.564) (-1.60,2.432) (-1.50,2.312) (-1.40,2.208) (-1.30,2.122) (-1.20,2.056) (-1.10,2.014) (-1.00,2.000) (-0.90,2.003) (-0.80,2.013) (-0.70,2.028) (-0.60,2.049) (-0.50,2.074) (-0.40,2.104) (-0.30,2.138) (-0.20,2.175) (-0.10,2.216) (0.00,2.259) (0.10,2.305) (0.20,2.352) (0.30,2.401) (0.40,2.450) (0.50,2.500) (0.60,2.550) (0.70,2.599) (0.80,2.648) (0.90,2.695) (1.00,2.741) (1.10,2.784) (1.20,2.825) (1.30,2.862) (1.40,2.896) (1.50,2.926) (1.60,2.951) (1.70,2.972) (1.80,2.987) (1.90,2.997) (2.00,3.000) (2.10,2.988) (2.20,2.952) (2.30,2.894) (2.40,2.816) (2.50,2.719) (2.60,2.604) (2.70,2.473) (2.80,2.328) (2.90,2.170) (3.00,2.000) (3.10,1.820) (3.20,1.632) (3.30,1.437) (3.40,1.236) (3.50,1.031) (3.60,0.824) (3.70,0.616) (3.80,0.408) (3.90,0.202) (4.00,0.000) (4.10,-0.202) (4.20,-0.408) (4.30,-0.616) (4.40,-0.824) (4.50,-1.031) (4.60,-1.236) (4.70,-1.437) (4.80,-1.632) (4.90,-1.820) (5.00,-2.000) (5.10,-2.170) (5.20,-2.328) (5.30,-2.473) (5.40,-2.604) (5.50,-2.719) (5.60,-2.816) (5.70,-2.894) (5.80,-2.952) (5.90,-2.988) (6.00,-3.000) (6.10,-2.962) (6.20,-2.856) (6.30,-2.694) (6.40,-2.488) (6.50,-2.250) (6.60,-1.992) (6.70,-1.726) (6.80,-1.464) (6.90,-1.218) (7.00,-1.000)};
    \node at (-6.3,3) {$\mathcal{C}_f$};
    #1
  \end{tikzpicture}%
}

% Exercice 3 : repère vide de -4 à 5
\newcommand{\repereVide}[1]{%
  \begin{tikzpicture}[x=0.75cm, y=0.75cm]
    \axes{-4}{5}{-4}{5}
    \foreach \x in {-4,...,-1,1,2,...,5} \node[below] at (\x,0) {\scriptsize $\x$};
    \foreach \y in {-4,...,-1,1,2,...,5} \node[left] at (0,\y) {\scriptsize $\y$};
    \node[below left] at (0,0) {\scriptsize $0$};
    #1
  \end{tikzpicture}%
}

% Exercice 5 : solde du compte (unité en ordonnée : 100 €)
\newcommand{\repereBanque}[1]{%
  \begin{tikzpicture}[x=1cm, y=0.6cm]
    \draw[gray!35, very thin, step=1] (0,-2) grid (12,6);
    \draw[-{Latex[length=2mm]}] (0,0) -- (12.5,0) node[right] {mois};
    \draw[-{Latex[length=2mm]}] (0,-2) -- (0,6.4) node[above] {solde (€)};
    \foreach \x in {1,...,12} \node[below] at (\x,0) {\scriptsize $\x$};
    \foreach \y/\l in {-2/-200,-1/-100,1/100,2/200,3/300,4/400,5/500,6/600} \node[left] at (0,\y) {\scriptsize $\l$};
    \node[below left] at (0,0) {\scriptsize $0$};
    \draw[very thick] plot coordinates {(0.00,1.000) (0.10,1.150) (0.20,1.299) (0.30,1.447) (0.40,1.592) (0.50,1.734) (0.60,1.873) (0.70,2.007) (0.80,2.136) (0.90,2.259) (1.00,2.375) (1.10,2.484) (1.20,2.584) (1.30,2.675) (1.40,2.757) (1.50,2.828) (1.60,2.888) (1.70,2.936) (1.80,2.971) (1.90,2.993) (2.00,3.000) (2.10,2.958) (2.20,2.840) (2.30,2.663) (2.40,2.440) (2.50,2.188) (2.60,1.920) (2.70,1.653) (2.80,1.400) (2.90,1.177) (3.00,1.000) (3.10,0.860) (3.20,0.736) (3.30,0.627) (3.40,0.528) (3.50,0.438) (3.60,0.352) (3.70,0.269) (3.80,0.184) (3.90,0.095) (4.00,0.000) (4.10,-0.109) (4.20,-0.232) (4.30,-0.363) (4.40,-0.496) (4.50,-0.625) (4.60,-0.744) (4.70,-0.847) (4.80,-0.928) (4.90,-0.981) (5.00,-1.000) (5.10,-0.981) (5.20,-0.928) (5.30,-0.847) (5.40,-0.744) (5.50,-0.625) (5.60,-0.496) (5.70,-0.363) (5.80,-0.232) (5.90,-0.109) (6.00,0.000) (6.10,0.100) (6.20,0.200) (6.30,0.300) (6.40,0.400) (6.50,0.500) (6.60,0.600) (6.70,0.700) (6.80,0.800) (6.90,0.900) (7.00,1.000) (7.10,1.098) (7.20,1.194) (7.30,1.287) (7.40,1.381) (7.50,1.475) (7.60,1.571) (7.70,1.671) (7.80,1.774) (7.90,1.884) (8.00,2.000) (8.10,2.130) (8.20,2.278) (8.30,2.442) (8.40,2.619) (8.50,2.806) (8.60,3.001) (8.70,3.200) (8.80,3.402) (8.90,3.602) (9.00,3.800) (9.10,3.992) (9.20,4.174) (9.30,4.346) (9.40,4.503) (9.50,4.644) (9.60,4.765) (9.70,4.864) (9.80,4.938) (9.90,4.984) (10.00,5.000) (10.10,4.983) (10.20,4.933) (10.30,4.853) (10.40,4.744) (10.50,4.609) (10.60,4.451) (10.70,4.271) (10.80,4.072) (10.90,3.856) (11.00,3.625) (11.10,3.382) (11.20,3.128) (11.30,2.866) (11.40,2.599) (11.50,2.328) (11.60,2.056) (11.70,1.785) (11.80,1.517) (11.90,1.255) (12.00,1.000)};
    #1
  \end{tikzpicture}%
}

% Exercice 7 : f(x) = -x^2 + 2x + 8 et g(x) = 3(x-3)^2 - 3 sur [-3 ; 5]
\newcommand{\repereFG}[1]{%
  \begin{tikzpicture}[x=0.9cm, y=0.4cm]
    \axes{-3}{5}{-7}{10}
    \foreach \x in {-3,...,-1,1,2,...,5} \node[below] at (\x,0) {\scriptsize $\x$};
    \foreach \y in {-6,-4,-2,2,4,...,10} \node[left] at (0,\y) {\scriptsize $\y$};
    \node[below left] at (0,0) {\scriptsize $0$};
    \begin{scope}
      \clip (-3,-7) rectangle (5,10);
      \draw[very thick, domain=-3:5, samples=80, smooth] plot (\x, {-\x*\x+2*\x+8});
      \draw[very thick, dashed, domain=1.1:5, samples=80, smooth] plot (\x, {3*(\x-3)*(\x-3)-3});
    \end{scope}
    \node at (-1.3,4.5) {$\mathcal{C}_f$};
    \node at (4.6,7) {$\mathcal{C}_g$};
    #1
  \end{tikzpicture}%
}

% Exercice 8 : tarifs de location (A : 5x ; B : 15 + 2x)
\newcommand{\repereVelo}[1]{%
  \begin{tikzpicture}[x=1.2cm, y=0.13cm]
    \draw[gray!20, very thin, xstep=0.5, ystep=2.5] (0,0) grid (10,50);
    \draw[gray!45, very thin, xstep=1, ystep=5] (0,0) grid (10,50);
    \draw[-{Latex[length=2mm]}] (0,0) -- (10.5,0) node[right] {durée (h)};
    \draw[-{Latex[length=2mm]}] (0,0) -- (0,53) node[above] {prix (€)};
    \foreach \x in {1,...,10} \node[below] at (\x,0) {\scriptsize $\x$};
    \foreach \y in {5,10,...,50} \node[left] at (0,\y) {\scriptsize $\y$};
    \node[below left] at (0,0) {\scriptsize $0$};
    \draw[very thick] (0,0) -- (10,50) node[pos=0.88, left] {tarif A};
    \draw[very thick, dashed] (0,15) -- (10,35) node[pos=0.9, below right] {tarif B};
    #1
  \end{tikzpicture}%
}

% Exercice 9 : triangle ABC (AB = 4, AC = 8) et rectangle AMNP
\newcommand{\triangleABC}{%
  \begin{tikzpicture}[scale=0.6]
    \fill[gray!25] (0,0) rectangle (5,1.5);
    \draw[thick] (0,0) node[below left] {$A$} -- (8,0) node[below right] {$C$} -- (0,4) node[above left] {$B$} -- cycle;
    \draw[thick] (0,1.5) node[left] {$M$} -- (5,1.5) node[above right] {$N$} -- (5,0) node[below] {$P$};
    \draw[{Latex[length=1.5mm]}-{Latex[length=1.5mm]}] (-1,0) -- node[left] {$x$} (-1,1.5);
  \end{tikzpicture}%
}
\newcommand{\repereRect}[1]{%
  \begin{tikzpicture}[x=2.6cm, y=0.5cm]
    \draw[gray!20, very thin, xstep=0.25, ystep=0.5] (0,0) grid (4,9);
    \draw[gray!45, very thin, xstep=0.5, ystep=1] (0,0) grid (4,9);
    \draw[-{Latex[length=2mm]}] (0,0) -- (4.2,0) node[right] {$x$};
    \draw[-{Latex[length=2mm]}] (0,0) -- (0,9.5) node[above] {aire (cm$^2$)};
    \foreach \x/\l in {0.5/0{,}5,1/1,1.5/1{,}5,2/2,2.5/2{,}5,3/3,3.5/3{,}5,4/4} \node[below] at (\x,0) {\scriptsize $\l$};
    \foreach \y in {1,...,9} \node[left] at (0,\y) {\scriptsize $\y$};
    \node[below left] at (0,0) {\scriptsize $0$};
    \draw[very thick, domain=0:4, samples=60, smooth] plot (\x, {\x*(8-2*\x)});
    #1
  \end{tikzpicture}%
}

% Exercice 10 : rectangle ABCD (AB = 6, AD = 14), x = 2
\newcommand{\rectABCD}{%
  \begin{tikzpicture}[scale=0.55]
    \fill[pattern=north east lines] (0,0) rectangle (2,2);
    \fill[pattern=north east lines] (2,2) rectangle (14,6);
    \draw[very thick] (0,0) node[below left] {$A$} rectangle (14,6);
    \node[above left] at (0,6) {$B$};
    \node[above right] at (14,6) {$C$};
    \node[below right] at (14,0) {$D$};
    \draw[thick] (0,2) node[left] {$M$} -- (2,2) node[above left] {$N$} -- (2,0) node[below] {$P$};
    \draw[thick] (2,2) -- (2,6) node[above] {$J$};
    \draw[thick] (2,2) -- (14,2) node[right] {$I$};
    \node[left] at (0,1) {$x$};
    \node[below] at (1,0) {$x$};
  \end{tikzpicture}%
}
\newcommand{\repereABCD}[1]{%
  \begin{tikzpicture}[x=2cm, y=0.07cm]
    \draw[gray!20, very thin, xstep=0.25, ystep=2] (0,0) grid (6,90);
    \draw[gray!45, very thin, xstep=0.5, ystep=10] (0,0) grid (6,90);
    \draw[-{Latex[length=2mm]}] (0,0) -- (6.3,0) node[right] {$x$};
    \draw[-{Latex[length=2mm]}] (0,0) -- (0,95);
    \foreach \x in {1,...,6} \node[below] at (\x,0) {\scriptsize $\x$};
    \foreach \y in {10,20,...,90} \node[left] at (0,\y) {\scriptsize $\y$};
    \node[below left] at (0,0) {\scriptsize $0$};
    \draw[very thick, domain=0:6, samples=60, smooth] plot (\x, {2*\x*\x-20*\x+84});
    \draw[very thick, domain=0:6, samples=60, smooth] plot (\x, {-2*\x*\x+20*\x});
    #1
  \end{tikzpicture}%
}

% Exercice 11 : bénéfice B(x) = -2x^2 + 26x - 60 sur [3 ; 10]
\newcommand{\repereBallon}[1]{%
  \begin{tikzpicture}[x=1.15cm, y=0.2cm]
    \draw[gray!20, very thin, xstep=0.5, ystep=1] (0,0) grid (11,30);
    \draw[gray!45, very thin, xstep=1, ystep=5] (0,0) grid (11,30);
    \draw[-{Latex[length=2mm]}] (0,0) -- (11.4,0) node[right] {prix (€)};
    \draw[-{Latex[length=2mm]}] (0,0) -- (0,31.5) node[above] {bénéfice (milliers d'€)};
    \foreach \x in {1,...,11} \node[below] at (\x,0) {\scriptsize $\x$};
    \foreach \y in {5,10,...,30} \node[left] at (0,\y) {\scriptsize $\y$};
    \node[below left] at (0,0) {\scriptsize $0$};
    \draw[very thick, domain=3:10, samples=60, smooth] plot (\x, {-2*\x*\x+26*\x-60});
    #1
  \end{tikzpicture}%
}

\begin{document}
\pagestyle{fancy}

\fancyhf{} % clear existing header/footer entries
\fancyhead[L]{ {\Niveau}}
\fancyhead[C]{ {\Chapitre}}
\fancyhead[R]{ {\Document}}
\fancyfoot[L]{ A. Fernandez }
\fancyfoot[R]{{\thepage\  / \pageref{LastPage}}}

% =====================================================
\begin{exercice}{1}
\begin{enumerate}[label=\arabic*.]
  \item $D_f = [-7\,;\,7]$.
  \item $f(-5) = 3$ \quad $f(2) = 3$ \quad $f(6) = -3$.
  \item Antécédents de $0$ : $-6$ et $4$. \quad Antécédent de $5$ : $-3$. \quad Antécédent de $-3$ : $6$.
  \item Non : la courbe ne descend jamais jusqu'à $y = -4$.
  \item
  \renewcommand{\arraystretch}{1.5}
  \begin{tabular}{|c|*{6}{C{1.3cm}|}}
    \hline
    $x$ & $-7$ & $-3$ & $-1$ & $7$ & $-5$ (ou $-2$ ou $2$) & $-6$ (ou $4$) \\ \hline
    $f(x)$ & $-2$ & $5$ & $2$ & $-1$ & $3$ & $0$ \\ \hline
  \end{tabular}
\end{enumerate}
\end{exercice}

% =====================================================
\begin{exercice}{2}
\begin{enumerate}[label=\arabic*.]
  \item \begin{enumerate}[label=\alph*.]
    \item $f(2) = 3 \times 2 - 5 = 1$ ; \quad $f(-1) = -3 - 5 = -8$ ; \quad $f(0) = -5$ ; \quad $f\left(\dfrac{1}{3}\right) = 1 - 5 = -4$.
    \item $3x - 5 = 7 \iff 3x = 12 \iff x = 4$ : l'antécédent de $7$ est $4$.

    $3x - 5 = 0 \iff x = \dfrac{5}{3}$ : l'antécédent de $0$ est $\dfrac{5}{3}$.
  \end{enumerate}
  \item \begin{enumerate}[label=\alph*.]
    \item $g(-3) = 9 - 4 = 5$ ; \quad $g(0) = -4$ ; \quad $g\left(\dfrac{1}{2}\right) = \dfrac{1}{4} - 4 = -\dfrac{15}{4}$.
    \item $x^2 - 4 = 0 \iff x^2 = 4 \iff x = 2$ ou $x = -2$. Les antécédents de $0$ sont $-2$ et $2$.
    \item $x^2 - 4 = -5 \iff x^2 = -1$. Un carré est toujours positif : $-5$ n'a pas d'antécédent par $g$.
  \end{enumerate}
\end{enumerate}
\end{exercice}

% =====================================================
\begin{exercice}{3}
\begin{enumerate}[label=\arabic*.]
  \item Par exemple : $f(-3) = 0{,}5 \times 9 - 0{,}5 \times (-3) - 3 = 4{,}5 + 1{,}5 - 3 = 3$.

  \renewcommand{\arraystretch}{1.5}
  \begin{tabular}{|c|*{8}{C{1.2cm}|}}
    \hline
    $x$ & $-3$ & $-2$ & $-1$ & $0$ & $1$ & $2$ & $3$ & $4$ \\ \hline
    $f(x)$ & $3$ & $0$ & $-2$ & $-3$ & $-3$ & $-2$ & $0$ & $3$ \\ \hline
  \end{tabular}
  \item
  \begin{center}
  \repereVide{
    \draw[very thick, blue!70!black, domain=-3.4:4.4, samples=60, smooth] plot (\x, {0.5*\x*\x-0.5*\x-3});
    \foreach \x in {-3,...,4} \fill[blue!70!black] (\x, {0.5*\x*\x-0.5*\x-3}) circle (2.5pt);
  }
  \end{center}
  \item $f(1) = -3 = y_A$, donc $A \in \mathcal{C}_f$.
  \item $f(5) = 0{,}5 \times 25 - 2{,}5 - 3 = 12{,}5 - 5{,}5 = 7 \neq 6$, donc $B \notin \mathcal{C}_f$.
\end{enumerate}
\end{exercice}

% =====================================================
\begin{exercice}{4}
\begin{enumerate}[label=\arabic*.]
  \item $f(2) = 8 - 6 + 1 = 3$ : $A \in \mathcal{C}_f$.

  $f(-1) = 2 + 3 + 1 = 6 \neq 5$ : $B \notin \mathcal{C}_f$.

  $f\left(\dfrac{1}{2}\right) = 2 \times \dfrac{1}{4} - \dfrac{3}{2} + 1 = \dfrac{1}{2} - \dfrac{3}{2} + 1 = 0$ : $C \in \mathcal{C}_f$.
  \item $f(-2) = 2 \times 4 + 6 + 1 = 15$ : $D(-2\,;\,15)$.
  \item $f(0) = 1$ : la courbe coupe l'axe des ordonnées en $(0\,;\,1)$.
\end{enumerate}
\end{exercice}

% =====================================================
\begin{exercice}{5}
\begin{enumerate}[label=\arabic*.]
  \item Le 1\textsuperscript{er} janvier : $100$~€. Au bout de 10 mois : $500$~€.
  \item Le solde vaut $100$~€ pour $x = 0$, $x = 3$, $x = 7$ et $x = 12$.
  \item Le solde est négatif pour $x \in\ ]4\,;\,6[$ (entre le 4\textsuperscript{e} et le 6\textsuperscript{e} mois).
  \item $x \in [0\,;\,3] \cup [7\,;\,12]$.
  \item Le solde le plus élevé est $500$~€, au bout de 10 mois.
\end{enumerate}
\end{exercice}

% =====================================================
\begin{exercice}{6}
\begin{enumerate}[label=\arabic*.]
  \item
  \begin{tasks}(4)
    \task $S = \{-6\,;\,4\}$
    \task $S = \{-5\,;\,-2\,;\,2\}$
    \task $S = \{-3\}$
    \task $S = \varnothing$
  \end{tasks}
  \item
  \begin{tasks}(2)
    \task $S = [-7\,;\,-6[\ \cup\ ]4\,;\,7]$
    \task $S = [-6\,;\,4]$
    \task $S = ]-5\,;\,-2[$
    \task $S = [-7\,;\,7]$ (la courbe ne dépasse jamais $5$)
  \end{tasks}
  \item $S = [-5\,;\,-2] \cup \{2\}$ : la courbe touche la droite $y = 3$ en $x = 2$ sans passer au-dessus. Le nombre $2$ est une solution isolée.
\end{enumerate}
\end{exercice}

% =====================================================
\begin{exercice}{7}
\begin{tasks}[label=\arabic*.](2)
  \task $f(x) = g(x)$ : $S = \{1\,;\,4\}$
  \task $f(x) > 8$ : $S = ]0\,;\,2[$
  \task $g(x) \leqslant 0$ : $S = [2\,;\,4]$
  \task $f(x) > g(x)$ : $S = ]1\,;\,4[$
  \task $f(x) \geqslant g(x)$ : $S = [1\,;\,4]$
  \task $f(x) < g(x)$ : $S = [-3\,;\,1[\ \cup\ ]4\,;\,5]$
  \task $f(x) \leqslant g(x)$ : $S = [-3\,;\,1] \cup [4\,;\,5]$
\end{tasks}
\end{exercice}

% =====================================================
\begin{exercice}{8}
\begin{enumerate}[label=\arabic*.]
  \item Tarif A : $20$~€. Tarif B : $23$~€.
  \item Les droites se coupent pour $x = 5$ (prix : $25$~€).
  \item La droite B est strictement en dessous de la droite A pour $x \in\ ]5\,;\,10]$.
  \item On trace la droite horizontale $y = 30$. Tarif A : $x \in [0\,;\,6]$, soit au plus 6 heures. Tarif B : $x \in [0\,;\,7{,}5]$, soit au plus 7 h 30.

  \begin{center}
  \repereVelo{
    \draw[red!70!black, dashed] (0,30) -- (10,30);
    \draw[red!70!black, densely dotted] (6,0) -- (6,30) (7.5,0) -- (7.5,30);
  }
  \end{center}
\end{enumerate}
\end{exercice}

% =====================================================
\begin{exercice}{9}
\begin{enumerate}[label=\arabic*.]
  \item $M$ est sur $[AB]$ et $AB = 4$, donc $D_f = [0\,;\,4]$.
  \item $f(x) = AM \times MN = x(8 - 2x)$.
  \item $f(1) = 1 \times 6 = 6$ ; \quad $f(1{,}5) = 1{,}5 \times 5 = 7{,}5$.
  \item \begin{enumerate}[label=\alph*.]
    \item $f(x) = 6$ pour $x = 1$ et $x = 3$.
    \item $f(x) \geqslant 6$ pour $x \in [1\,;\,3]$.
    \item L'aire semble maximale pour $x = 2$ (milieu de $[AB]$) : elle vaut alors $8$~cm$^2$.
  \end{enumerate}
  \item $(MN) \parallel (AC)$ (toutes deux perpendiculaires à $(AB)$). D'après le théorème de Thalès dans le triangle $ABC$ :
  $\dfrac{MN}{AC} = \dfrac{BM}{BA}$, soit $\dfrac{MN}{8} = \dfrac{4 - x}{4}$, donc $MN = 2(4 - x) = 8 - 2x$.
\end{enumerate}
\end{exercice}

% =====================================================
\begin{exercice}{10}
\begin{enumerate}[label=\arabic*.]
  \item $M$ est sur $[AB]$ et $AB = 6$ : $x \in [0\,;\,6]$.
  \item \begin{enumerate}[label=\alph*.]
    \item $NI = 14 - x$ et $NJ = 6 - x$. Donc
    $f(x) = x^2 + (14 - x)(6 - x) = x^2 + 84 - 14x - 6x + x^2 = 2x^2 - 20x + 84$.
    \item L'aire de $ABCD$ est $6 \times 14 = 84$, donc $g(x) = 84 - f(x) = 84 - 2x^2 + 20x - 84 = -2x^2 + 20x$.
  \end{enumerate}
  \item $f(0) = 84$ et $g(0) = 0$ : la courbe de $f$ est celle qui part de $84$, la courbe de $g$ part de $0$.
  \item Les courbes se coupent pour $x = 3$. Vérification : $f(3) = 18 - 60 + 84 = 42$ et $g(3) = -18 + 60 = 42$. Il faut placer $M$ à $3$ unités de $A$, au milieu de $[AB]$.
  \item L'aire hachurée semble la plus petite pour $x = 5$ : elle vaut alors $f(5) = 50 - 100 + 84 = 34$.
\end{enumerate}
\end{exercice}

% =====================================================
\begin{exercice}{11}
\textbf{Partie A.}
\begin{enumerate}[label=\arabic*.]
  \item $f(5) = 20 - 10 = 10$ : $10\,000$ ballons. \quad $f(9) = 20 - 18 = 2$ : $2\,000$ ballons.
  \item $20 - 2x = 6 \iff 2x = 14 \iff x = 7$ : prix de $7$~€. \quad $20 - 2x = 12 \iff x = 4$ : prix de $4$~€.
\end{enumerate}

\textbf{Partie B.}
\begin{enumerate}[label=\arabic*.]
  \item \begin{enumerate}[label=\alph*.]
    \item Sur chaque ballon, l'entreprise gagne $x - 3$ euros, et elle vend $20 - 2x$ milliers de ballons. Donc $B(x) = (20 - 2x)(x - 3)$ milliers d'euros.
    \item $B(x) = 20x - 60 - 2x^2 + 6x = -2x^2 + 26x - 60$.
  \end{enumerate}
  \item \begin{enumerate}[label=\alph*.]
    \item $B(6) = 8 \times 3 = 24$ : $24\,000$~€ de bénéfice.
    \item $B(2) = 16 \times (-1) = -16$ : l'entreprise perd $16\,000$~€ (le prix est inférieur au coût de fabrication).
  \end{enumerate}
  \item
  \renewcommand{\arraystretch}{1.5}
  \begin{tabular}{|c|*{9}{C{1cm}|}}
    \hline
    $x$ & $1$ & $2$ & $3$ & $4$ & $5$ & $6$ & $7$ & $8$ & $9$ \\ \hline
    $B(x)$ & $-36$ & $-16$ & $0$ & $12$ & $20$ & $24$ & $24$ & $20$ & $12$ \\ \hline
  \end{tabular}
  \item \begin{enumerate}[label=\alph*.]
    \item $B(x) \geqslant 20$ pour $x \in [5\,;\,8]$.
    \item $B(x) < 12$ pour $x \in [3\,;\,4[\ \cup\ ]9\,;\,10]$.
    \item Le bénéfice est le plus grand pour un prix de $6{,}50$~€ : il vaut alors $24{,}5$ milliers d'euros, soit $24\,500$~€.
  \end{enumerate}
\end{enumerate}
\end{exercice}
\end{document}
